Los experimentos fueron ejecutados en una máquina bajo Linux (kernel 6.x) equipada con un procesador \textbf{AMD Ryzen 5 7600X} (6 núcleos, 12 hilos, 4.7~GHz base / 5.3~GHz boost), tarjeta gráfica \textbf{AMD Radeon RX 6900 XT} (16~GB VRAM) y \textbf{32~GB de memoria RAM DDR5} en dual-channel ($2 \times 16\,\text{GB}$). El código se desarrolló en C++17 y fue compilado mediante \texttt{g++} con optimización \texttt{-O2}. Los tiempos de ejecución se registraron con \texttt{std::chrono::high\_resolution\_clock}.

Para ordenamiento, se evaluaron arreglos con tamaños $N \in \{10^1, 10^3, 10^5, 10^7\}$ bajo cuatro escenarios: secuencias aleatorias, ordenadas ascendentemente, ordenadas descendentemente y arreglos con dominios restringidos de valores. Para la multiplicación matricial, se utilizaron matrices cuadradas con dimensiones $N \in \{16, 64, 256, 1024\}$, variando su estructura interna (dispersas, diagonales y densas) y dominio numérico ($D_0 \in \{0, 1\}$ y $D_{10} \in [0, 10]$), con tres réplicas independientes por configuración.